% --- Archon physics formalization source begin ---
% archon:physics
% archon:covers PhyXMiniProblems/problem_phyx_mini_0797.lean
% archon:source-report reports/phyx_mini/problem_phyx_mini_0797.source.json

\chapter{Physics problem phyx\_mini\_0797}
\label{ch:PhyXMiniProblems_problem_phyx_mini_0797}

\paragraph{Problem source.}
\#\# Physical scenario

An airplane's compass indicates that it is headed due north, and its air-speed indicator shows that it is moving through the air at \textbackslash{}( 240 \textbackslash{}, \textbackslash{}text\{km/h\} \textbackslash{}) as shown in figure. If there is a \textbackslash{}( 100 \textbackslash{}, \textbackslash{}text\{km/h\} \textbackslash{}) wind from west to east.

\#\# Question

what is the velocity of the airplane relative to the earth?

\#\# Answer choices

- A: A: \textbackslash{}( 160 \textbackslash{}

- B: B: \textbackslash{}( 260 \textbackslash{}

- C: C: \textbackslash{}( 290 \textbackslash{}

- D: D: \textbackslash{}( 410 \textbackslash{}

\#\# Figure caption (auxiliary; use the image as primary evidence)

The image is a vector diagram involving an airplane and its velocities. There are three airplanes depicted, showing movement at different points. The vectors are represented as arrows indicating direction and magnitude.

1. **Vectors:**
   - **\textbackslash{}[ \textbackslash{}vec\{v\}\_\{A/E\} \textbackslash{}]** (Air relative to Earth): A horizontal vector pointing to the right, labeled "100 km/h, east."
   - **\textbackslash{}[ \textbackslash{}vec\{v\}\_\{P/A\} \textbackslash{}]** (Plane relative to Air): A vertical vector pointing upwards, labeled "240 km/h, north."
   - **\textbackslash{}[ \textbackslash{}vec\{v\}\_\{P/E\} \textbackslash{}]** (Plane relative to Earth): A diagonal vector formed by the sum of the other two vectors. 

2. **Planes:**
   - Three planes are illustrated to show the relative motion along each vector path, positioned at the starting, intermediate, and end points.

3. **Angles and Directions:**
   - An angle \textbackslash{}( \textbackslash{}alpha \textbackslash{}) is indicated between the north direction and the resultant velocity vector \textbackslash{}[ \textbackslash{}vec\{v\}\_\{P/E\} \textbackslash{}].
   - A compass rose is shown with directions labeled N, E, S, and W for orientation.

This setup likely represents the problem of finding the resultant velocity of the plane relative to the Earth, taking into account the wind's effect coming from the east.

\#\# Dataset metadata

Kinematics; Spatial Relation Reasoning, Multi-Formula Reasoning

\paragraph{Recorded answer/context.}
B: B: \textbackslash{}( 260 \textbackslash{}

\paragraph{Figure/image path.}
/root/proposal\_for\_physic/science-mango/phyx\_mini\_run/phyx\_data/test\_image/797.png

\paragraph{Formalization target.}
create a compiling Lean file with sorry bodies at `PhyXMiniProblems/problem\_phyx\_mini\_0797.lean`. The Lean declarations must preserve the physical quantities, dimensions or dimensional roles, figure labels, governing-law hypotheses, and final relation expressed by this problem.
Use Mathlib/Physlib names found through LeanExplore where available. If a physics API is missing, introduce faithful local abstractions rather than scalar placeholder aliases.

\begin{theorem}[Physics formalization target]
\label{thm:physics:phyx_mini_0797:target}
\lean{PhyXMiniProblems.ProblemPhyXMini0797.problem_phyx_mini_0797}
\uses{def:physics:phyx-mini-0797:phyxminiproblems-problemphyxmini0797-velocityinkilometersperhour, def:physics:phyx-mini-0797:phyxminiproblems-problemphyxmini0797-speedinkilometersperhour, def:physics:phyx-mini-0797:phyxminiproblems-problemphyxmini0797-easthat, def:physics:phyx-mini-0797:phyxminiproblems-problemphyxmini0797-northhat, def:physics:phyx-mini-0797:phyxminiproblems-problemphyxmini0797-airplanewindsetup, def:physics:phyx-mini-0797:phyxminiproblems-problemphyxmini0797-matchesstateddirections, def:physics:phyx-mini-0797:phyxminiproblems-problemphyxmini0797-matchesproblemreadouts, def:physics:phyx-mini-0797:phyxminiproblems-problemphyxmini0797-matchessuppliedfigure, def:physics:phyx-mini-0797:phyxminiproblems-problemphyxmini0797-obeysdirectionalmeasurementmodel, def:physics:phyx-mini-0797:phyxminiproblems-problemphyxmini0797-obeysgalileanvelocitycomposition, def:physics:phyx-mini-0797:phyxminiproblems-problemphyxmini0797-representsfiguredriftangle, def:physics:phyx-mini-0797:phyxminiproblems-problemphyxmini0797-matchesdisplayedgroundspeed}
The assigned autoformalize agent should translate this physics problem into Lean declarations in the covered file, with theorem and lemma proof bodies written as `by sorry`.
\end{theorem}
\begin{proof}
This is an autoformalization task, not a proof task. Produce statements that can later be proved by the physics prover without weakening the physical model.
\end{proof}

% --- Archon declaration topology begin ---
\section{Declaration topology}
\begin{definition}[velocity Dimension]
  \label{def:physics:phyx-mini-0797:phyxminiproblems-problemphyxmini0797-velocitydimension}
  \lean{PhyXMiniProblems.ProblemPhyXMini0797.velocityDimension}
  The physical dimension L T⁻¹ of velocity
\end{definition}
\begin{proof}
  This is the declared model definition of velocity Dimension.
\end{proof}
\begin{definition}[Planar Velocity Quantity]
  \label{def:physics:phyx-mini-0797:phyxminiproblems-problemphyxmini0797-planarvelocityquantity}
  \lean{PhyXMiniProblems.ProblemPhyXMini0797.PlanarVelocityQuantity}
  \uses{def:physics:phyx-mini-0797:phyxminiproblems-problemphyxmini0797-velocitydimension}
  A physical planar velocity, independent of the units chosen to read its two signed Cartesian components
\end{definition}
\begin{proof}
  This is the declared model definition of Planar Velocity Quantity.
\end{proof}
\begin{definition}[kilometer Hour Units]
  \label{def:physics:phyx-mini-0797:phyxminiproblems-problemphyxmini0797-kilometerhourunits}
  \lean{PhyXMiniProblems.ProblemPhyXMini0797.kilometerHourUnits}
  A coherent unit choice whose length and time units are kilometres and hours
\end{definition}
\begin{proof}
  This is the declared model definition of kilometer Hour Units.
\end{proof}
\begin{definition}[velocity Readout]
  \label{def:physics:phyx-mini-0797:phyxminiproblems-problemphyxmini0797-velocityreadout}
  \lean{PhyXMiniProblems.ProblemPhyXMini0797.velocityReadout}
  \uses{def:physics:phyx-mini-0797:phyxminiproblems-problemphyxmini0797-planarvelocityquantity}
  Read a physical planar velocity in an arbitrary coherent unit system
\end{definition}
\begin{proof}
  This is the declared model definition of velocity Readout.
\end{proof}
\begin{definition}[velocity In Kilometers Per Hour]
  \label{def:physics:phyx-mini-0797:phyxminiproblems-problemphyxmini0797-velocityinkilometersperhour}
  \lean{PhyXMiniProblems.ProblemPhyXMini0797.velocityInKilometersPerHour}
  \uses{def:physics:phyx-mini-0797:phyxminiproblems-problemphyxmini0797-planarvelocityquantity, def:physics:phyx-mini-0797:phyxminiproblems-problemphyxmini0797-kilometerhourunits, def:physics:phyx-mini-0797:phyxminiproblems-problemphyxmini0797-velocityreadout}
  Read the signed Cartesian components of a velocity in kilometres per hour
\end{definition}
\begin{proof}
  This is the declared model definition of velocity In Kilometers Per Hour.
\end{proof}
\begin{definition}[speed In Kilometers Per Hour]
  \label{def:physics:phyx-mini-0797:phyxminiproblems-problemphyxmini0797-speedinkilometersperhour}
  \lean{PhyXMiniProblems.ProblemPhyXMini0797.speedInKilometersPerHour}
  \uses{def:physics:phyx-mini-0797:phyxminiproblems-problemphyxmini0797-planarvelocityquantity, def:physics:phyx-mini-0797:phyxminiproblems-problemphyxmini0797-velocityinkilometersperhour}
  The speed (Euclidean norm) of a velocity, in kilometres per hour
\end{definition}
\begin{proof}
  This is the declared model definition of speed In Kilometers Per Hour.
\end{proof}
\begin{definition}[east Hat]
  \label{def:physics:phyx-mini-0797:phyxminiproblems-problemphyxmini0797-easthat}
  \lean{PhyXMiniProblems.ProblemPhyXMini0797.eastHat}
  The dimensionless unit vector pointing east in the diagram
\end{definition}
\begin{proof}
  This is the declared model definition of east Hat.
\end{proof}
\begin{definition}[north Hat]
  \label{def:physics:phyx-mini-0797:phyxminiproblems-problemphyxmini0797-northhat}
  \lean{PhyXMiniProblems.ProblemPhyXMini0797.northHat}
  The dimensionless unit vector pointing north in the diagram
\end{definition}
\begin{proof}
  This is the declared model definition of north Hat.
\end{proof}
\begin{definition}[Physical System]
  \label{def:physics:phyx-mini-0797:phyxminiproblems-problemphyxmini0797-physicalsystem}
  \lean{PhyXMiniProblems.ProblemPhyXMini0797.PhysicalSystem}
  The three systems occurring in the relative-velocity subscripts
\end{definition}
\begin{proof}
  This is the declared model definition of Physical System.
\end{proof}
\begin{definition}[Cardinal Direction]
  \label{def:physics:phyx-mini-0797:phyxminiproblems-problemphyxmini0797-cardinaldirection}
  \lean{PhyXMiniProblems.ProblemPhyXMini0797.CardinalDirection}
  Cardinal headings printed by the compass rose or stated in the prose
\end{definition}
\begin{proof}
  This is the declared model definition of Cardinal Direction.
\end{proof}
\begin{definition}[Diagram Direction]
  \label{def:physics:phyx-mini-0797:phyxminiproblems-problemphyxmini0797-diagramdirection}
  \lean{PhyXMiniProblems.ProblemPhyXMini0797.DiagramDirection}
  Directions used for the three arrows in the supplied vector diagram
\end{definition}
\begin{proof}
  This is the declared model definition of Diagram Direction.
\end{proof}
\begin{definition}[Relative Velocity Label]
  \label{def:physics:phyx-mini-0797:phyxminiproblems-problemphyxmini0797-relativevelocitylabel}
  \lean{PhyXMiniProblems.ProblemPhyXMini0797.RelativeVelocityLabel}
  The three literal relative-velocity labels in image 797.png
\end{definition}
\begin{proof}
  This is the declared model definition of Relative Velocity Label.
\end{proof}
\begin{definition}[moving System]
  \label{def:physics:phyx-mini-0797:phyxminiproblems-problemphyxmini0797-movingsystem}
  \lean{PhyXMiniProblems.ProblemPhyXMini0797.movingSystem}
  \uses{def:physics:phyx-mini-0797:phyxminiproblems-problemphyxmini0797-physicalsystem, def:physics:phyx-mini-0797:phyxminiproblems-problemphyxmini0797-relativevelocitylabel}
  The moving system named by the numerator of a velocity label
\end{definition}
\begin{proof}
  This is the declared model definition of moving System.
\end{proof}
\begin{definition}[reference System]
  \label{def:physics:phyx-mini-0797:phyxminiproblems-problemphyxmini0797-referencesystem}
  \lean{PhyXMiniProblems.ProblemPhyXMini0797.referenceSystem}
  \uses{def:physics:phyx-mini-0797:phyxminiproblems-problemphyxmini0797-physicalsystem, def:physics:phyx-mini-0797:phyxminiproblems-problemphyxmini0797-relativevelocitylabel}
  The reference system named by the denominator of a velocity label
\end{definition}
\begin{proof}
  This is the declared model definition of reference System.
\end{proof}
\begin{definition}[cardinal Unit Vector]
  \label{def:physics:phyx-mini-0797:phyxminiproblems-problemphyxmini0797-cardinalunitvector}
  \lean{PhyXMiniProblems.ProblemPhyXMini0797.cardinalUnitVector}
  \uses{def:physics:phyx-mini-0797:phyxminiproblems-problemphyxmini0797-easthat, def:physics:phyx-mini-0797:phyxminiproblems-problemphyxmini0797-northhat, def:physics:phyx-mini-0797:phyxminiproblems-problemphyxmini0797-cardinaldirection}
  The unit vector assigned to each compass direction
\end{definition}
\begin{proof}
  This is the declared model definition of cardinal Unit Vector.
\end{proof}
\begin{definition}[Diagram Point]
  \label{def:physics:phyx-mini-0797:phyxminiproblems-problemphyxmini0797-diagrampoint}
  \lean{PhyXMiniProblems.ProblemPhyXMini0797.DiagramPoint}
  The three vertices of the head-to-tail velocity triangle
\end{definition}
\begin{proof}
  This is the declared model definition of Diagram Point.
\end{proof}
\begin{definition}[Airplane Velocity Figure]
  \label{def:physics:phyx-mini-0797:phyxminiproblems-problemphyxmini0797-airplanevelocityfigure}
  \lean{PhyXMiniProblems.ProblemPhyXMini0797.AirplaneVelocityFigure}
  \uses{def:physics:phyx-mini-0797:phyxminiproblems-problemphyxmini0797-cardinaldirection, def:physics:phyx-mini-0797:phyxminiproblems-problemphyxmini0797-diagramdirection, def:physics:phyx-mini-0797:phyxminiproblems-problemphyxmini0797-relativevelocitylabel, def:physics:phyx-mini-0797:phyxminiproblems-problemphyxmini0797-diagrampoint}
  Typed evidence exposed by the supplied raster. A missing magnitude is represented by none; in particular, the image does not print the requested magnitude beside v\_P/E
\end{definition}
\begin{proof}
  This is the declared model definition of Airplane Velocity Figure.
\end{proof}
\begin{definition}[Airplane Wind Setup]
  \label{def:physics:phyx-mini-0797:phyxminiproblems-problemphyxmini0797-airplanewindsetup}
  \lean{PhyXMiniProblems.ProblemPhyXMini0797.AirplaneWindSetup}
  \uses{def:physics:phyx-mini-0797:phyxminiproblems-problemphyxmini0797-planarvelocityquantity, def:physics:phyx-mini-0797:phyxminiproblems-problemphyxmini0797-cardinaldirection, def:physics:phyx-mini-0797:phyxminiproblems-problemphyxmini0797-relativevelocitylabel, def:physics:phyx-mini-0797:phyxminiproblems-problemphyxmini0797-airplanevelocityfigure}
  The independent physical velocities, instrument/report readouts, and diagram. The airplane-earth velocity is a field to be constrained by the composition law; it is not defined from answer choice B or from the number 260
\end{definition}
\begin{proof}
  This is the declared model definition of Airplane Wind Setup.
\end{proof}
\begin{definition}[Matches Stated Directions]
  \label{def:physics:phyx-mini-0797:phyxminiproblems-problemphyxmini0797-matchesstateddirections}
  \lean{PhyXMiniProblems.ProblemPhyXMini0797.MatchesStatedDirections}
  \uses{def:physics:phyx-mini-0797:phyxminiproblems-problemphyxmini0797-airplanewindsetup}
  The two directions stated in the prose
\end{definition}
\begin{proof}
  This is the declared model definition of Matches Stated Directions.
\end{proof}
\begin{definition}[Matches Problem Readouts]
  \label{def:physics:phyx-mini-0797:phyxminiproblems-problemphyxmini0797-matchesproblemreadouts}
  \lean{PhyXMiniProblems.ProblemPhyXMini0797.MatchesProblemReadouts}
  \uses{def:physics:phyx-mini-0797:phyxminiproblems-problemphyxmini0797-airplanewindsetup}
  The numerical instrument and weather-report readouts from the problem
\end{definition}
\begin{proof}
  This is the declared model definition of Matches Problem Readouts.
\end{proof}
\begin{definition}[Matches Supplied Figure]
  \label{def:physics:phyx-mini-0797:phyxminiproblems-problemphyxmini0797-matchessuppliedfigure}
  \lean{PhyXMiniProblems.ProblemPhyXMini0797.MatchesSuppliedFigure}
  \uses{def:physics:phyx-mini-0797:phyxminiproblems-problemphyxmini0797-cardinaldirection, def:physics:phyx-mini-0797:phyxminiproblems-problemphyxmini0797-airplanewindsetup}
  The labels, directions, endpoints, magnitudes, airplanes, compass rose, and angle mark visible in image 797.png. No ground-speed or numerical alpha answer is included here
\end{definition}
\begin{proof}
  This is the declared model definition of Matches Supplied Figure.
\end{proof}
\begin{definition}[Obeys Directional Measurement Model]
  \label{def:physics:phyx-mini-0797:phyxminiproblems-problemphyxmini0797-obeysdirectionalmeasurementmodel}
  \lean{PhyXMiniProblems.ProblemPhyXMini0797.ObeysDirectionalMeasurementModel}
  \uses{def:physics:phyx-mini-0797:phyxminiproblems-problemphyxmini0797-velocityinkilometersperhour, def:physics:phyx-mini-0797:phyxminiproblems-problemphyxmini0797-cardinalunitvector, def:physics:phyx-mini-0797:phyxminiproblems-problemphyxmini0797-airplanewindsetup}
  The measurement interpretation used in the textbook model: the compass and air-speed indicator determine the plane-through-air vector, while the wind report determines the air-through-earth vector. This law contains neither the airplane-earth vector nor its requested speed
\end{definition}
\begin{proof}
  This is the declared model definition of Obeys Directional Measurement Model.
\end{proof}
\begin{definition}[Obeys Galilean Velocity Composition]
  \label{def:physics:phyx-mini-0797:phyxminiproblems-problemphyxmini0797-obeysgalileanvelocitycomposition}
  \lean{PhyXMiniProblems.ProblemPhyXMini0797.ObeysGalileanVelocityComposition}
  \uses{def:physics:phyx-mini-0797:phyxminiproblems-problemphyxmini0797-velocityreadout, def:physics:phyx-mini-0797:phyxminiproblems-problemphyxmini0797-airplanewindsetup}
  The Galilean relative-velocity composition law, stated in every coherent unit system rather than only for the numerical kilometre/hour readout
\end{definition}
\begin{proof}
  This is the declared model definition of Obeys Galilean Velocity Composition.
\end{proof}
\begin{definition}[Represents Figure Drift Angle]
  \label{def:physics:phyx-mini-0797:phyxminiproblems-problemphyxmini0797-representsfiguredriftangle}
  \lean{PhyXMiniProblems.ProblemPhyXMini0797.RepresentsFigureDriftAngle}
  \uses{def:physics:phyx-mini-0797:phyxminiproblems-problemphyxmini0797-velocityinkilometersperhour, def:physics:phyx-mini-0797:phyxminiproblems-problemphyxmini0797-northhat, def:physics:phyx-mini-0797:phyxminiproblems-problemphyxmini0797-airplanewindsetup}
  The figure-derived role of alpha: it is the undirected angle from the north unit vector to the airplane-earth velocity vector. The raster supplies the role and symbol, but no numerical value
\end{definition}
\begin{proof}
  This is the declared model definition of Represents Figure Drift Angle.
\end{proof}
\begin{definition}[Answer Choice]
  \label{def:physics:phyx-mini-0797:phyxminiproblems-problemphyxmini0797-answerchoice}
  \lean{PhyXMiniProblems.ProblemPhyXMini0797.AnswerChoice}
  Labels of the four multiple-choice answers
\end{definition}
\begin{proof}
  This is the declared model definition of Answer Choice.
\end{proof}
\begin{definition}[displayed Ground Speed Kilometers Per Hour]
  \label{def:physics:phyx-mini-0797:phyxminiproblems-problemphyxmini0797-displayedgroundspeedkilometersperhour}
  \lean{PhyXMiniProblems.ProblemPhyXMini0797.displayedGroundSpeedKilometersPerHour}
  \uses{def:physics:phyx-mini-0797:phyxminiproblems-problemphyxmini0797-answerchoice}
  The speed in kilometres per hour printed beside each answer label
\end{definition}
\begin{proof}
  This is the declared model definition of displayed Ground Speed Kilometers Per Hour.
\end{proof}
\begin{definition}[Matches Displayed Ground Speed]
  \label{def:physics:phyx-mini-0797:phyxminiproblems-problemphyxmini0797-matchesdisplayedgroundspeed}
  \lean{PhyXMiniProblems.ProblemPhyXMini0797.MatchesDisplayedGroundSpeed}
  \uses{def:physics:phyx-mini-0797:phyxminiproblems-problemphyxmini0797-speedinkilometersperhour, def:physics:phyx-mini-0797:phyxminiproblems-problemphyxmini0797-airplanewindsetup, def:physics:phyx-mini-0797:phyxminiproblems-problemphyxmini0797-answerchoice, def:physics:phyx-mini-0797:phyxminiproblems-problemphyxmini0797-displayedgroundspeedkilometersperhour}
  A displayed choice agrees with the physical airplane-earth speed
\end{definition}
\begin{proof}
  This is the declared model definition of Matches Displayed Ground Speed.
\end{proof}
\begin{definition}[recorded Dataset Answer]
  \label{def:physics:phyx-mini-0797:phyxminiproblems-problemphyxmini0797-recordeddatasetanswer}
  \lean{PhyXMiniProblems.ProblemPhyXMini0797.recordedDatasetAnswer}
  \uses{def:physics:phyx-mini-0797:phyxminiproblems-problemphyxmini0797-answerchoice}
  The answer label recorded in the dataset metadata
\end{definition}
\begin{proof}
  This is the declared model definition of recorded Dataset Answer.
\end{proof}
% --- Archon declaration topology end ---
% --- Archon physics formalization source end ---
